\documentclass[aspectratio=169]{beamer}
% Shared Beamer configuration for MUS2640 lecture decks (included after \documentclass).
\usepackage[utf8]{inputenc}
\usepackage[T1]{fontenc}
\usepackage{lmodern}
\usepackage{graphicx}
\usepackage{booktabs}

\usetheme{Madrid}
\usecolortheme{dolphin}
\usefonttheme{professionalfonts}

\setbeamercovered{transparent}
\setbeamertemplate{navigation symbols}{}
\setbeamertemplate{footline}[frame number]

\newcommand{\coursename}{MUS2640 · Sensing Sound and Music}
\newcommand{\instituteuo}{University of Oslo · Department of Musicology / RITMO}


\title{Week 3 · Acoustics}
\subtitle{\coursename}
\institute{\instituteuo}
\date{}

\begin{document}

\begin{frame}
  \titlepage
\end{frame}

\begin{frame}{Aims}
  \begin{itemize}
    \item Chain \textbf{cause $\rightarrow$ propagation $\rightarrow$ reception}.
    \item Basic wave concepts: frequency, amplitude, spectrum; periodic vs aperiodic.
    \item Rooms and instruments as resonant systems (qualitative + demos).
  \end{itemize}
\end{frame}

\begin{frame}{Roadmap}
  \begin{enumerate}
    \item Nature of sound waves; vibrations.
    \item Spectral pictures (connect to listening vocabulary).
    \item Room acoustics: reflections, reverberation, critical distance (ideas).
    \item Instrument families — excitation and radiation (overview).
  \end{enumerate}
\end{frame}

\begin{frame}{Cause and effect}
  \begin{itemize}
    \item Source mechanisms determine partial structure (harmonic vs noisy).
    \item Space filters and mixes arrivals — affects timbre and rhythm clarity.
    \item Same score / recording can ``sound different'' in another room.
  \end{itemize}
\end{frame}

\begin{frame}{Why acoustics for musicologists?}
  \begin{itemize}
    \item Bridges verbal description and measurable quantities.
    \item Explains limits of recording/reproduction (preview: electroacoustics).
    \item Supports critical listening to performance venues and mixes.
  \end{itemize}
\end{frame}

\begin{frame}{Optional activity}
  \begin{itemize}
    \item Sonic Visualiser / notebook: load a short clip; inspect waveform + spectrogram.
    \item Relate a spectral feature to an adjective you would use in Week 2 language.
  \end{itemize}
\end{frame}

\begin{frame}{Outside class}
  \begin{itemize}
    \item Work chapter exercises; sketch a concept map from source $\rightarrow$ ear.
    \item Clap in two contrasting rooms; estimate $T_{60}$ from a 20\,dB decay and compare with your ears.
  \end{itemize}
\end{frame}

\begin{frame}{Summary}
  \begin{itemize}
    \item Physical acoustics grounds later psychoacoustics and studio chains.
    \item Next: \textbf{Psychoacoustics} — ear and brain mapping physics to experience.
  \end{itemize}
\end{frame}

\end{document}
